% Part: first-order-logic
% Chapter: syntax-and-semantics
% Section: formation-sequences

\documentclass[../../../include/open-logic-section]{subfiles}

\begin{document}

\olfileid{pl}{syn}{fseq}

\olsection{构造序列}

通过归纳定义来定义公式，并通过归纳法证明公式的性质，是一种优雅而高效的方法。然而，考虑一种更自底向上、逐步构造公式的方法也会很有用，我们在此借助\emph{构造序列}的概念来实现这一点。

\begin{defn}[公式的构造序列]
\ollabel{defn:fseq-frm}
语言~$\Lang L_0$ 中的符号串组成的一个有限序列$\tuple{!A_0,\dotsc,!A_n}$是 $!A$ 的\emph{构造序列}，如果 $!A \ident !A_n$，并且对所有 $i \leq n$，要么 $!A_i$ 是原子公式，要么存在 $j,k < i$ 使得以下之一成立：
\begin{enumerate}
    \tagitem{prvNot}{$!A_i \ident \lnot !A_j$。}{}%
    \tagitem{prvAnd}{$!A_i \ident (!A_j \land !A_k)$。}{}%
    \tagitem{prvOr}{$!A_i \ident (!A_j \lor !A_k)$。}{}%
    \tagitem{prvIf}{$!A_i \ident (!A_j \lif !A_k)$。}{}%
    \tagitem{prvIff}{$!A_i \ident (!A_j \liff !A_k)$。}{}%
\end{enumerate}
\end{defn}

\begin{ex}
\[
\tuple{
    \Obj p_0,
    \Obj p_1,
    (\Obj p_1 \land \Obj p_0),
    \lnot (\Obj p_1 \land \Obj p_0)
}
\]
是
$\lnot (\Obj p_1 \land \Obj p_0)$
的一个构造序列，
\[
\tuple{
    \Obj p_0,
    \Obj p_1,
    \Obj p_0,
    (\Obj p_1 \land \Obj p_0),
    (\Obj p_0 \lif \Obj p_1),
    \lnot (\Obj p_1 \land \Obj p_0)
}.
\]
也是。%
从第二个例子可以看出，构造序列可能包含“垃圾”：冗余或对构造没有贡献的公式。
\end{ex}

\begin{prop}
\ollabel{prop:formed}
$\Frm[L_0]$ 中的每个公式~$!A$ 都有一个构造序列。
\end{prop}

\begin{proof}
假设 $!A$ 是原子公式。那么序列~$\tuple{!A}$ 就是~$!A$ 的一个构造序列。
%
现在假设 $!B$ 和~$!C$ 分别有构造序列
$\tuple{!B_0,\dotsc,!B_n}$ 和 $\tuple{!C_0,\dotsc,!C_m}$。
%
\begin{enumerate}
  \tagitem{prvNot}{若 $!A \ident \lnot !B$，
      则 $\tuple{!B_0,\dotsc,!B_n,\lnot !B_n}$
      是~$!A$ 的一个构造序列。}{}
  \tagitem{prvAnd}{若 $!A \ident (!B \land !C)$，
      则 $\tuple{!B_0,\dotsc,!B_n,!C_0,\dotsc,!C_m,(!B_n \land !C_m)}$
      是~$!A$ 的一个构造序列。}{}
  \tagitem{prvOr}{若 $!A \ident (!B \lor !C)$，
      则 $\tuple{!B_0,\dotsc,!B_n,!C_0,\dotsc,!C_m,(!B_n \lor !C_m)}$
      是~$!A$ 的一个构造序列。}{}
  \tagitem{prvIf}{若 $!A \ident (!B \lif !C)$，
      则 $\tuple{!B_0,\dotsc,!B_n,!C_0,\dotsc,!C_m,(!B_n \lif !C_m)}$
      是~$!A$ 的一个构造序列。}{}
  \tagitem{prvIff}{若 $!A \ident (!B \liff !C)$，
      则 $\tuple{!B_0,\dotsc,!B_n,!C_0,\dotsc,!C_m,(!B_n \liff !C_m)}$
      是~$!A$ 的一个构造序列。}{}
\end{enumerate}
由公式上的归纳原理，每个公式都有一个构造序列。
\end{proof}

我们也可以证明其逆命题。这一点很重要，因为它表明我们定义公式的两种方式是等价的：它们给出相同的结果。这也意味着，我们可以通过对构造序列的长度做普通归纳来证明关于公式的定理。

\begin{lem}
\ollabel{lem:fseq-init}
设 $\tuple{!A_0,\dotsc,!A_n}$ 是~$!A_n$ 的一个构造序列，且 $k \leq n$。则 $\tuple{!A_0,\dotsc,!A_k}$
是~$!A_k$ 的一个构造序列。
\end{lem}

\begin{proof}
练习。
\end{proof}

\begin{thm}
\ollabel{thm:fseq-frm-equiv}
$\Frm[L_0]$ 是语言~$\Lang L_0$ 中所有具有构造序列的符号串的集合。
\end{thm}

\begin{proof}
令 $F$ 为语言~$\Lang L_0$ 中所有具有构造序列的符号串的集合。
我们已在 \olref[pl][syn][fseq]{prop:formed} 中看到
$\Frm[L_0] \subseteq F$，现在证明反向包含。

假设 $!A$ 有一个构造序列 $\tuple{!A_0,\dotsc,!A_n}$。
我们通过对~$n$ 进行强归纳来证明 $!A \in \Frm[L_0]$。
归纳假设为：每个具有长度 $m < n$ 的构造序列的符号串都属于~$\Frm[L_0]$。
根据构造序列的定义，$!A_n$ 要么是原子公式，要么存在 $j,k < n$ 使得以下之一成立：
\begin{enumerate}
    \tagitem{prvNot}{$!A_n \ident \lnot !A_j$。}{}%
    \tagitem{prvAnd}{$!A_n \ident (!A_j \land !A_k)$。}{}%
    \tagitem{prvOr}{$!A_n \ident (!A_j \lor !A_k)$。}{}%
    \tagitem{prvIf}{$!A_n \ident (!A_j \lif !A_k)$。}{}%
    \tagitem{prvIff}{$!A_n \ident (!A_j \liff !A_k)$。}{}%
\end{enumerate}
现在我们分情况讨论。如果 $!A_n$ 是原子公式，那么 $!A_n \in \Frm[L_0]$。
否则，假设 $!A \equiv
(!A_j \land !A_k)$。根据 \olref[pl][syn][fseq]{lem:fseq-init},
$\tuple{!A_0,\dotsc,!A_j}$ 和 $\tuple{!A_0,\dotsc,!A_k}$ 分别是
$!A_j$ 和~$!A_k$ 的构造序列。由于它们是~$!A$ 的构造序列的真初始子序列，它们的长度都小于~$n$。
因此根据归纳假设，$!A_j$ 和~$!A_k$ 都属于~$\Frm[L_0]$，
于是根据公式的定义，$(!A_j \land !A_k)$ 也属于~$\Frm[L_0]$。其他情况可类似推得。
\end{proof}

\end{document}